\documentclass[a4paper]{article}
% Español (México) con XeLaTeX: fontspec para fuentes Unicode, polyglossia
% para la división silábica y los nombres del documento en la variante mexicana.
\usepackage{fontspec}
\usepackage{polyglossia}
\setdefaultlanguage[variant=mexican]{spanish}
% Los nombres chinos van como «pinyin (original)»: xeCJK da a los caracteres
% Han su propia fuente; sin ella XeLaTeX los omite con un aviso.
% FandolSong no cubre algunos caracteres de nombres (U+73FA); WenQuanYi Zen Hei sí.
\usepackage[AutoFallBack=true]{xeCJK}
\setCJKmainfont{FandolSong-Regular.otf}
\setCJKfallbackfamilyfont{\CJKrmdefault}{WenQuanYi Zen Hei}
% polyglossia no traduce los nombres de listings.
\AtBeginDocument{\providecommand{\lstlistingname}{}\renewcommand{\lstlistingname}{Listado}%
  \providecommand{\lstlistlistingname}{}\renewcommand{\lstlistlistingname}{Índice de listados}}
\usepackage{amsmath, amssymb}
\usepackage{graphicx}
\usepackage[margin=2.35cm]{geometry}
\usepackage[most]{tcolorbox}
\usepackage{booktabs}
\usepackage{tabularx}
\usepackage{longtable}
\usepackage{array}
\usepackage{float}
\usepackage{xcolor}
\usepackage{hyperref}
\usepackage{fancyhdr}
\hypersetup{colorlinks=true, linkcolor=blue!55!black, urlcolor=blue!60!black}
\graphicspath{{./}{figures/}}
\setlength{\parskip}{0.55em}
\setlength{\parindent}{2em}
\setlength{\emergencystretch}{3em}
\renewcommand{\arraystretch}{1.18}
\newtcolorbox{knowledgebox}[1]{enhanced, colback=blue!5!white, colframe=blue!70!black, colbacktitle=blue!70!black, coltitle=white, fonttitle=\bfseries, title={#1}, sharp corners, breakable}
\newtcolorbox{importantbox}[1]{enhanced, colback=yellow!10!white, colframe=yellow!75!black, colbacktitle=yellow!75!black, coltitle=black, fonttitle=\bfseries, title={#1}, sharp corners, breakable}
\newtcolorbox{warningbox}[1]{enhanced, colback=red!5!white, colframe=red!70!black, colbacktitle=red!70!black, coltitle=white, fonttitle=\bfseries, title={#1}, sharp corners, breakable}
\pagestyle{fancy}
\fancyhf{}
\fancyfoot[C]{\thepage}
\fancyfoot[R]{\footnotesize\color{gray}\href{https://github.com/hqhq1025/ai-course-notes}{AI Course Notes}}
\renewcommand{\headrulewidth}{0pt}
\renewcommand{\footrulewidth}{0pt}
\newcommand{\videourl}{https://www.youtube.com/watch?v=n4_c_HsodPg}
\newcommand{\repourl}{https://github.com/hqhq1025/ai-course-notes}

\begin{document}
\begin{titlepage}
\centering
{\Large Notas didácticas de una entrevista a fondo\par}
\vspace{1.0cm}
{\huge\bfseries IA del mundo físico: de Waymo y Momenta a Galaxea (星海图)\par}
\vspace{0.5cm}
{\Large Zhang Xiaojun conversa con Gao Jiyang (高继扬)\par}
\vspace{0.35cm}
{\large Elaborado a partir del video público de Zhang Xiaojun Podcast, la transcripción hecha con GPU, la estructura de capítulos y diagramas conceptuales propios\par}
\vspace{0.3cm}
{\large 2026-02-13\par}
\vspace{0.85cm}
\includegraphics[width=0.88\textwidth,height=0.42\textheight,keepaspectratio]{cover.jpg}\par
\vfill
\begin{tcolorbox}[width=0.92\textwidth, colback=black!2!white, colframe=black!55, sharp corners]
\textbf{Título del video}: 132. Entrevista de 3 horas con Gao Jiyang, fundador de Galaxea: el bagre, Zeng Guofan (曾国藩), las dos caras de Waymo y Momenta, un lobo y la salida de Xu Huazhe (许华哲)\par
\textbf{Canal del video}: Zhang Xiaojun Podcast\par
\textbf{Duración del video}: 03:04:52\par
\textbf{Enlace del video}: \href{\videourl}{\nolinkurl{\videourl}}\par
\textbf{Descripción de los materiales}: YouTube no tiene subtítulos disponibles; el texto principal se elaboró a partir de la transcripción hecha con faster-whisper large-v3 en CUDA, los capítulos del video, la portada y figuras didácticas propias. Las tomas fijas de la entrevista no se repiten en el texto principal.\par
\textbf{Página del proyecto}: \href{\repourl}{\nolinkurl{\repourl}}\par
\end{tcolorbox}
\end{titlepage}

\tableofcontents
\newpage
\section{Introducción: por qué emprender en robótica «no es romántico»}

Esta sección fija primero la manera de leer el episodio: no conviene tomarlo como una noticia de chismes sobre «quién dejó qué empresa», ni solo como publicidad corporativa de Xinghaitu (星海图). Una lectura más valiosa es verlo como el testimonio oral de cómo un emprendedor traslada su formación en conducción autónoma a la inteligencia corporizada.

En apariencia, esta entrevista trata de la trayectoria personal de Gao Jiyang (高继扬), fundador de Xinghaitu; en realidad es una clase sobre la industria acerca de «cómo crece una empresa de IA del mundo físico». Al inicio, el conductor plantea una duda: ¿por qué en el sector chino de la inteligencia corporizada casi no aparecen personas con el marcado romanticismo tecnológico que se ve en los emprendimientos de modelos grandes? La respuesta de Gao Jiyang atraviesa toda la conversación: la cadena de la robótica es larguísima, los ciclos son larguísimos y el equipo tiene, por naturaleza, que meter la cabeza en la tierra.

Aquí «en la tierra» no es una expresión negativa, sino las restricciones reales del mundo físico: el equipo completo, la cadena de suministro, la recolección de datos, las instalaciones del cliente, la latencia en el dispositivo, el mantenimiento, los canales de distribución, el financiamiento y el crecimiento de la organización. Los modelos grandes o las aplicaciones de software pueden arrancar a partir del modelo y de la difusión, pero a una empresa de robótica le resulta muy difícil esquivar el hardware y el valor para el cliente. Por eso, estas notas organizan el episodio como un «campo de entrenamiento industrial de la IA del mundo físico»: Waymo le enseñó ingeniería de sistemas, Momenta le enseñó la entrega en producción en masa, y Xinghaitu recombina el equipo completo, los datos, los modelos y el ciclo cerrado con el cliente.

\begin{importantbox}{Tesis central del episodio}
La inteligencia corporizada no es «una empresa de algoritmos con una carcasa de robot», sino una cadena larga que va del equipo completo, la cadena de suministro, el sistema de datos, la AI infra y los modelos algorítmicos hasta los canales de distribución y el valor para el cliente. A corto plazo los algoritmos pueden difundirse con rapidez; a largo plazo, la barrera competitiva tiene más probabilidades de surgir del ciclo cerrado con el mundo real.
\end{importantbox}

\begin{knowledgebox}{Nota sobre la estrategia visual}
Este video es una toma fija de entrevista, sin slides, pizarrón ni demostraciones de producto. Conforme al estándar para notas de pódcast, el cuerpo del texto no repite fotogramas de las personas; la portada sirve para identificar la fuente, y en el cuerpo se usan diagramas conceptuales, como el modelo de negocio de la conducción autónoma, el volante de datos y el sistema dual del cerebro robótico, para explicar el contenido de la entrevista.
\end{knowledgebox}

\subsection{Resumen de la sección}

Este episodio no es una simple historia de fundador, sino un caso de cómo la experiencia en conducción autónoma se traslada a la inteligencia corporizada. Más adelante se abordarán, en orden, la metodología de Gao Jiyang, la comparación entre Waymo y Momenta, las decisiones de Xinghaitu sobre el equipo completo y la cadena de suministro, la Data Recipe, el cerebro robótico VLM/VLA y las disyuntivas organizacionales detrás de «la salida de Xu Huazhe (许华哲)».
\section{De la síntesis inductiva a la IA del mundo físico}

Gao Jiyang (高继扬) vuelve una y otra vez a un mismo método: inducir y sintetizar. La primaria, la olimpiada de física, las publicaciones del doctorado, las entrevistas de trabajo y la decisión de emprender las aborda de la misma manera: «primero fijar el objetivo, luego descomponer el camino y después aumentar la probabilidad de acertar». El método parece sencillo, pero explica por qué su trayectoria lo llevó del Departamento de Ingeniería Electrónica de la Universidad Tsinghua (清华), pasando por el aprendizaje profundo, la visión por computadora y la conducción autónoma, hasta la inteligencia corpórea.

\begin{figure}[H]
\centering
\includegraphics[width=0.96\textwidth]{trajectory-industrial-training.png}
\caption{La trayectoria de formación industrial de Gao Jiyang: síntesis inductiva, iniciación en la IA, Waymo, Momenta y Galaxea (星海图).}
\end{figure}

\begin{knowledgebox}{Lectura de la figura: cómo se acumulan tres capacidades}
La ruta de la figura no es una lista de cargos. La olimpiada de física y el doctorado le enseñaron a inducir, descomponer y gestionar probabilidades; las prácticas en SenseTime (商汤) le permitieron sentir por primera vez que una red neuronal puede extraer regularidades de los datos por sí sola; Waymo le dio formación en ingeniería de sistemas y en medición; Momenta, en la entrega de productos y el valor para el cliente; y Galaxea lleva todas estas capacidades a la inteligencia corpórea.
\end{knowledgebox}

\subsection{Zeng Guofan: del letrado purista a los logros prácticos}

La sección anterior trató la «síntesis inductiva» de Gao Jiyang; esta explica de dónde viene su realismo. El ejemplo de Zeng Guofan (曾国藩) no es un adorno histórico, sino uno de los primeros modelos con los que entendió la organización de una empresa emergente.

En la entrevista, Gao Jiyang cuenta que leyó a Zeng Guofan después de un tropiezo al solicitar ingreso a las universidades. No le interesaba la curiosidad histórica, sino «cómo un letrado confuciano purista se convierte en alguien capaz de movilizar recursos, organizar equipos y librar batallas difíciles en el mundo real». Este pasaje es fundamental para entender su visión del emprendimiento: para hacer cosas en el mundo real no bastan las ideas, los artículos ni la inteligencia; hay que movilizar recursos, organizar personas y asumir los resultados.

\begin{importantbox}{La perspectiva de los logros prácticos}
Los logros prácticos no son utilitarismo de corto plazo, sino convertir un objetivo en recursos, organización, acciones y resultados. Un emprendimiento de robótica necesita especialmente esta perspectiva, porque cada uno de sus pasos tiene que materializarse en el mundo real.
\end{importantbox}

\subsection{El atractivo del aprendizaje profundo}

Durante sus prácticas en SenseTime entrenó una red neuronal por primera vez y percibió la fascinación de que «la máquina extraiga regularidades de los datos por sí misma». En la programación tradicional, una persona resume las reglas y luego las escribe como código; en el aprendizaje automático, las reglas se condensan en los parámetros y el modelo aprende a partir de la distribución de los datos. Ese momento fijó su manera de entender la IA en adelante: el valor de la IA no está solo en automatizar, sino en cambiar la forma en que los seres humanos descubren regularidades y construyen productividad.

\begin{knowledgebox}{Términos clave: de las reglas al aprendizaje}
\begin{tabularx}{\textwidth}{p{0.22\textwidth}p{0.32\textwidth}X}
\toprule
Término & Problema que resuelve & Significado en este episodio \\
\midrule
Rule-based & Reglas, módulos e if-else escritos a mano por personas & Base importante de los sistemas tradicionales de conducción autónoma y de robótica. \\
Data-driven & Aprender regularidades a partir de los datos & La razón central por la que Gao Jiyang eligió la IA y la conducción autónoma. \\
Benchmark & Medir el desempeño de los modelos con una evaluación unificada & Lo necesitan la publicación de artículos, la iteración de modelos y la recipe de datos. \\
Physical World AI & IA que actúa directamente sobre tareas del mundo físico & La conducción autónoma es su primera forma; la inteligencia corpórea, la siguiente. \\
\bottomrule
\end{tabularx}
\end{knowledgebox}

\subsection{Resumen de la sección}

El hilo conductor de Gao Jiyang es convertir el esfuerzo personal en un método reutilizable y luego llevar la capacidad de la IA basada en datos al mundo físico. Lo que le atrae de la conducción autónoma y de los robots es precisamente que la IA se vuelve una variable de base de toda la industria, no una herramienta de optimización local.
\section{Waymo: ingeniería de sistemas y mentalidad de ingeniero}

Al terminar el doctorado eligió la conducción autónoma, porque en ese momento era la forma más clara de physical world AI. Lo que más aprendió en Waymo no fueron solo los algoritmos de conducción autónoma, sino cómo funciona el sistema completo: cómo se integran la percepción, la localización, los mapas de alta precisión fuera de línea, la predicción, la toma de decisiones, la planificación, el control, la simulación y la cadena de herramientas en la nube hasta formar un sistema confiable.

\subsection{La continuidad histórica de la arquitectura de conducción autónoma}

Al releer los artículos de conducción autónoma de alrededor de 2008, encontró que la arquitectura general de hacia 2018 no difería en lo fundamental de los primeros marcos: el sistema seguía dividido en módulos de percepción, localización, mapas, planificación y control. El verdadero cambio fue que, cuando las técnicas de IA maduraron, muchos módulos de percepción pasaron de clustering, reglas y métodos tradicionales a redes neuronales.

De ahí surge una distinción importante: la lógica de fondo de la gran arquitectura tradicional de conducción autónoma se parece más a la de robotics y da prioridad a la modularidad, la interpretabilidad y los corner case; la vía AI native da prioridad al enfoque basado en datos, al diseño de extremo a extremo y a la mejora del benchmark global. No se trata de decidir cuál tiene razón, sino de equilibrar de manera distinta la confiabilidad, la capacidad de iteración, la interpretabilidad y la escalabilidad.

\begin{warningbox}{Lo modular no es atraso, ni lo de extremo a extremo lo resuelve todo}
En un sistema modular es más fácil localizar los corner case y los límites de responsabilidad, pero puede tardar en cambiar de rumbo y acumular reglas locales; un sistema de extremo a extremo mejora con más facilidad las métricas globales a partir de datos, pero puede introducir fallas que no se pueden explicar. La IA del mundo físico a menudo necesita volver a trazar la frontera entre ambos.
\end{warningbox}

\subsection{Mentalidad de ingeniero: descomponer y medir}

Gao Jiyang (高继扬) resume que la formación central que le dio Waymo fue la mentalidad de ingeniero: descomponer y medir. Se divide un problema complejo en subproblemas manejables, y estos a su vez en código y pruebas; al mismo tiempo, se construye un sistema de medición que va de las métricas de nivel superior a las intermedias y hasta las pruebas unitarias de base. Esto difiere de resolver problemas en las olimpiadas de física, donde se trata sobre todo de reconocer a qué tipo de problema corresponde cada uno; un sistema de ingeniería, en cambio, tiene que funcionar de manera continua, localizar problemas de manera continua y volver de manera continua al objetivo global.

\begin{importantbox}{Mentalidad de ingeniero}
La mentalidad de ingeniero no consiste en «saber escribir código», sino en descomponer un sistema complejo en módulos verificables y conectar, mediante un sistema de métricas, los cambios locales con los resultados de nivel superior.
\end{importantbox}

\subsection{Resumen de la sección}

Waymo le dio a Gao Jiyang una visión del sistema completo y una formación en ingeniería: saber cómo se descompone un sistema de IA del mundo físico, cómo se mide y cómo se ejecuta a largo plazo. Pero estaba lejos de los clientes, del producto y de la gestión de una empresa, así que su siguiente paso fue ir a Momenta, más cerca de la entrega de producción en serie.
\section{Momenta: producción en serie, valor para el cliente y el «bagre»}

La sección sobre Waymo explica cómo se descompone un sistema; la sección sobre Momenta explica cómo los clientes ponen a prueba un sistema. Para Gao Jiyang (高继扬), pasar de Waymo a Momenta no fue saltar de una buena empresa a otra buena empresa, sino pasar del campo de entrenamiento de ingenieros al terreno de la entrega en producción en serie.

Si Waymo es el paraíso de los ingenieros, Momenta es el campo de batalla de la producción en serie. Gao Jiyang eligió Momenta porque quería pasar de los sistemas de conducción autónoma al producto, a los clientes y a la gestión de una empresa; en particular, quería entender cómo un sistema de nivel demo se convierte en un producto que pueda entregarse a un fabricante de automóviles.

\begin{figure}[H]
\centering
\includegraphics[width=0.96\textwidth]{waymo-momenta-contrast.png}
\caption{Waymo y Momenta: dos formas extremas de formación organizacional.}
\end{figure}

\begin{knowledgebox}{Lectura de la figura: qué forma cada organización}
A la izquierda, Waymo forma en ingeniería de sistemas, infra, calidad del código y ejecución de la estrategia a largo plazo; a la derecha, Momenta forma en entrega, comunicación con el cliente, ajuste organizacional y el ciclo de valor comercial. Una enseña cómo funciona un sistema; la otra enseña cómo los clientes usan un sistema, pagan por él y obligan a iterarlo.
\end{knowledgebox}

\subsection{Por qué la producción en serie es una ruta de datos}

La lógica central de Momenta es esta: la conducción autónoma necesita, en última instancia, una gran cantidad de datos reales. Si se depende solo de una flotilla propia, la cobertura de ciudades y escenarios avanza demasiado despacio; si capacidades como la asistencia a la conducción o el estacionamiento se instalan en vehículos de producción en serie, primero se crea valor para el cliente y después se obtienen datos a través de esos vehículos, con lo que se forma un ciclo de datos impulsado por el negocio. Lo decisivo aquí no es «venderle a los fabricantes de automóviles» en sí, sino que la obtención de datos deja de ser un costo puro y pasa a pagarse con el valor que recibe el cliente.

\begin{importantbox}{Ciclo de datos y valor comercial}
La producción en serie no es solo la entrega de un proyecto: hace que el producto llegue a vehículos reales, clientes reales y caminos reales, y así convierte la recolección de datos en una actividad impulsada por el negocio. Cuanto más claro es el valor para el cliente, más sostenible es el retorno de datos.
\end{importantbox}

\subsection{El papel de bagre}

Lo anterior trató la transformación organizacional de Momenta; aquí el enfoque vuelve al propio Gao Jiyang. Lo que se llama bagre no es una etiqueta de personalidad, sino una forma de intervenir en sistemas que se forjó una y otra vez en un entorno de entrega bajo alta presión.

El «bagre» del título de la entrevista se refiere a que en Momenta a Gao Jiyang lo asignaban con flexibilidad a distintos módulos: percepción, localización, estacionamiento, infra, planificación y control, y la producción en serie de NOA. Su propio resumen es: entrar rápido en áreas desconocidas, entender el sistema con una metodología fija, descomponer las tareas, asignar a las personas adecuadas, monitorear la retroalimentación, ampliar cuando la retroalimentación es buena y reducir y ajustar cuando no lo es.

\begin{knowledgebox}{La metodología del bagre}
El bagre no se dedica a agitarlo todo, sino a que el sistema reaccione mejor. Sus acciones incluyen: entrar en un área desconocida, trazar rápidamente un mapa, descomponer tareas, asignar personas, fijar indicadores, observar la retroalimentación y ajustar la organización.
\end{knowledgebox}

\subsection{Resumen de la sección}

Lo que Momenta le dio a Gao Jiyang fue formación en valor para el cliente y en presión organizacional. Le mostró que, para que un equipo técnico se convierta en una empresa de producto, tiene que pasar por la prueba de la entrega, la eliminación, el ajuste y la producción en serie.
\section{Modelos de negocio de la conducción autónoma: quién es dueño de los datos y de los clientes}

Después de exponer la trayectoria personal de entrenamiento en Waymo y Momenta, hay que pasar al nivel del modelo de negocio. Al igual que en la robótica, en la conducción autónoma la ruta tecnológica nunca es una elección puramente técnica: está ligada al punto de acceso a los datos, a la relación con los clientes y a la bolsa de utilidades.

En la entrevista, Gao Jiyang (高继扬) divide los modelos de negocio de la conducción autónoma en varias categorías: el robotaxi al estilo de Waymo, la venta de vehículos más suscripción de software de los fabricantes de automóviles, el proveedor al estilo de Momenta y la plataforma de utilidades del vehículo completo, similar a la de Huawei. Esta clasificación ayuda a entender por qué se bifurcan las rutas tecnológicas: quién es dueño del vehículo, quién es dueño del usuario, quién es dueño de los datos y quién asume la responsabilidad del servicio influye directamente en la arquitectura técnica y en las capacidades de la organización.

\begin{figure}[H]
\centering
\includegraphics[width=0.94\textwidth]{autonomous-driving-business-models.png}
\caption{Los cuatro modelos de negocio de la conducción autónoma: Robotaxi, suscripción del fabricante de automóviles, proveedor y plataforma de utilidades del vehículo completo.}
\end{figure}

\begin{knowledgebox}{Lectura de la figura: el modelo de negocio determina el ciclo de los datos}
El Robotaxi opera su propia flota: los datos y el servicio están en sus manos, pero su escalamiento es lento. La suscripción del fabricante de automóviles posee los vehículos y los usuarios, y obtiene los datos directamente. El proveedor necesita los proyectos de los fabricantes para llegar a la producción en serie, y tanto el valor para el cliente como el retorno de los datos dependen de la colaboración. La plataforma de utilidades del vehículo completo, por su parte, redefine la bolsa de utilidades del vehículo mediante la marca, los canales de venta, la conducción inteligente y la cabina inteligente.
\end{knowledgebox}

\subsection{Por qué esto es instructivo para la robótica}

La clasificación anterior proviene de la conducción autónoma, pero esta sección busca trasladarla: la robótica no cuenta con un mercado ya establecido como el automotriz ni con datos de flota de origen natural. Precisamente porque carece de estas condiciones previas, la experiencia de la conducción autónoma solo puede ofrecer un marco, no respuestas directas.

La inteligencia corpórea todavía no tiene un lado de la demanda ni una escala de flota tan maduros como los del automóvil. Lo que la conducción autónoma enseña a la robótica no es copiar el modelo de robotaxi o el de proveedor, sino ver con claridad que el ciclo cerrado de datos debe estar ligado al valor para el cliente. Si un robot no tiene escenarios reales ni unidades entregadas, no hay datos continuos; sin datos, es difícil entrenar un modelo fundacional de acciones; sin la capacidad del modelo, no se puede aumentar el valor para el cliente.

\begin{warningbox}{La experiencia de la conducción autónoma no puede copiarse mecánicamente}
El automóvil ya tiene un mercado enorme: los vehículos se venden. El cuerpo del robot de propósito general todavía no cuenta con una demanda igual de madura. La robótica debe crear al mismo tiempo la demanda de hardware, el valor de los escenarios y el ciclo cerrado de datos, lo cual es más difícil que «trasladar la ruta de la conducción autónoma».
\end{warningbox}

\subsection{Resumen de la sección}

El problema de fondo de los modelos de negocio de la conducción autónoma es: de dónde provienen los datos, cómo se forma el valor para el cliente y cómo itera el sistema. La inteligencia corpórea tampoco puede eludir estas preguntas; la diferencia es que su demanda, su hardware y sus escenarios están más dispersos.
\section{De un mal BP a la cadena de suministro del robot completo}

Ahora entramos en el robot de Xinghaitu (星海图). Las secciones anteriores explicaron por qué Gao Jiyang (高继扬) cree en la physical world AI y qué aprendió de la conducción autónoma; esta sección responde a la decisión más importante que tomó al fundar la empresa: por qué una empresa de inteligencia corporizada no puede dedicarse solo a los algoritmos.

El primer giro de Xinghaitu surgió de una conclusión: si se quiere hacer inteligencia corporizada, no basta con hacer la «inteligencia»; hay que fabricar el robot completo. Gao Jiyang dijo en la entrevista que, para que un robot cierre de verdad el ciclo entre datos y modelo, la empresa necesita su propio cuerpo robótico, su hardware, su cadena de suministro y sus escenarios de cliente. De lo contrario, por bueno que sea el modelo base, le falta una vía de acceso estable al mundo real.

\begin{figure}[H]
\centering
\includegraphics[width=0.96\textwidth]{three-product-combo.png}
\caption{Los tres componentes de Xinghaitu: robot completo, modelo base y herramientas de posentrenamiento.}
\end{figure}

\begin{knowledgebox}{Lectura de la figura: por qué tres componentes}
Los tres componentes de la figura apuntan a una misma experiencia objetivo: capacitar a un robot como se capacita a un empleado. El robot completo aporta la vía de acceso física y la capacidad de recopilar datos; el modelo base aporta la base general de acciones y de comprensión; las herramientas de posentrenamiento permiten que el cliente adapte el robot a su escenario con pocas demostraciones y con práctica autónoma. Si falta cualquiera de los tres, es muy difícil llevar el robot de la demo a escenarios productivos.
\end{knowledgebox}

\subsection{El ritmo de 2024, 2025 y 2026}

Los tres componentes describen la estructura final, pero una startup no puede desplegar todas sus capacidades a la vez. Esta subsección se centra en el ritmo, porque el ritmo determina si la organización puede soportar la complejidad.

El ritmo que plantea Gao Jiyang es el siguiente: en 2024 el foco está en el robot completo y la cadena de suministro; en 2025, en la inteligencia de datos; en 2026 empieza el trabajo en la cadena de suministro de largo plazo. Este ritmo refleja la idea de «avanzar paso a paso, consolidando cada posición»: se mantiene la estrategia de largo plazo, pero no se despliega todo al mismo tiempo en un mismo periodo. Si una empresa de robótica persigue demasiado pronto y en paralelo el hardware, los modelos, las aplicaciones, los canales y los clientes, la complejidad termina por desbordar a la organización.

\begin{importantbox}{Implicaciones estratégicas de avanzar paso a paso}
Avanzar paso a paso no es ser conservador, sino avanzar según las dependencias. Sin el robot completo y la cadena de suministro no hay datos estables; sin datos estables no hay un buen VLA; sin un modelo utilizable es difícil generar valor para el cliente.
\end{importantbox}

\subsection{Resumen de la sección}

Xinghaitu eligió fabricar el robot completo no porque el hardware en sí sea más atractivo, sino porque el robot completo es la vía de acceso al ciclo cerrado entre datos, modelo y cliente. Para que una empresa de inteligencia corporizada lleve su inteligencia a fondo, primero tiene que consolidar la cadena física.
\section{Data Recipe: los datos no son un inventario, sino un ciclo cerrado}

Solo cuando ya se tienen el robot completo y la cadena de suministro tiene sentido hablar de datos. Este orden no puede invertirse: los datos de robots no son texto descargado de internet, sino algo que se produce en el hardware, los escenarios, las tareas y la retroalimentación de los clientes.

El «panorama de datos» de EP132 y el de EP134 se aclaran mutuamente. Lo central del Data Recipe del que habla aquí Gao Jiyang (高继扬) no es vender los datos como producto, sino conectar el robot completo, los escenarios, las tareas, la recolección, el entrenamiento, la evaluación y el ecosistema de código abierto en un volante iterable.

\begin{figure}[H]
\centering
\includegraphics[width=0.96\textwidth]{robot-data-recipe-flywheel.png}
\caption{Volante del Data Recipe para robots: despliegue del robot completo, recolección de datos, entrenamiento del modelo, valor en campo, retorno de errores y apertura compartida.}
\end{figure}

\begin{knowledgebox}{Lectura de la figura: por qué el Data Recipe no consiste en vender datos}
Los datos de la figura provienen de hardware y escenarios reales. Cuando el robot completo llega a escenarios de desarrolladores o de productividad, genera trayectorias, fallas y retroalimentación; estos datos entrenan VLA, modelos base y herramientas de posentrenamiento; cuando el modelo crea valor para el cliente, trae consigo más escenarios y más retroalimentación. Los datos no son un inventario estático, sino el combustible de aprendizaje que generan de manera continua el producto y los clientes.
\end{knowledgebox}

\subsection{Por qué importa llegar a diez mil unidades entregadas}

Gao Jiyang considera que «lograr entregas de diez mil unidades en escenarios de productividad» es una bet importante en este momento. Esta meta no es una simple cifra de ventas, sino una validación integral de los datos, la cadena de suministro, la calidad y el valor para el cliente. Diez mil unidades significan que el hardware puede entregarse, que los clientes están dispuestos a usarlo, que el sistema de mantenimiento puede seguirle el paso y que el retorno de datos puede escalarse.

\begin{knowledgebox}{Términos clave: el ciclo cerrado de datos de robots}
\begin{tabularx}{\textwidth}{p{0.23\textwidth}p{0.32\textwidth}X}
\toprule
Término & Problema que resuelve & Significado en este episodio \\
\midrule
Data Recipe & Receta de origen, selección, entrenamiento y retroalimentación de los datos & Determina cómo los datos reales se convierten en capacidades del modelo. \\
Demo in Video & Demostración que parece funcionar en un video & Tiende a generar expectativas demasiado altas. \\
Demo in Office & Demostración en vivo en la oficina de la empresa & Más sólida que el video, pero sigue siendo un entorno controlado. \\
Demo in the Wild & Demostración desplegada con clientes, en ferias y en escenarios de distintos países & Más cercana a la capacidad real de generalización. \\
Escenario de productividad & Los clientes usan el robot para realizar trabajo real & Fuente del ciclo cerrado entre datos y valor comercial. \\
\bottomrule
\end{tabularx}
\end{knowledgebox}

\subsection{Código abierto y datos compartidos}

Xinghaitu (星海图) publica en código abierto conjuntos de datos y modelos, y comparte datos con sus clientes. Gao Jiyang subraya que los datos no son su negocio principal, pero que con gusto los comparte. Esto puede entenderse como una estrategia de ecosistema: permitir que más desarrolladores y clientes validen tareas, hagan visibles sus necesidades y generen retroalimentación, para así ampliar los escenarios de uso de los modelos base y de la cadena de herramientas.

\subsection{Resumen de la sección}

La esencia del Data Recipe es conectar el robot completo, los clientes y el entrenamiento de modelos. Lo que de verdad escasea en una empresa de robótica no es «tener un lote de datos», sino un mecanismo que produzca de forma continua datos, evaluaciones y retroalimentación de alto valor.
\section{El cerebro del robot: el sistema dual de VLM y VLA}

El ciclo cerrado de datos resuelve «de qué se alimenta el modelo»; el cerebro del robot resuelve «cómo actúa el modelo». Aquí es también donde este episodio pasa del análisis de la industria a la estructura técnica.

Cuando el presentador pregunta «cómo se construye el cerebro», Gao Jiyang (高继扬) lo divide en dos modelos fundacionales: uno es el VLM de nivel superior, que se encarga de descomponer instrucciones, razonar y planificar tareas; el otro es el modelo fundacional de acción, el VLA, que a partir de vision y language produce action y hace que el cuerpo del robot ejecute la tarea. Esta es la estructura técnica más importante del episodio.

\begin{figure}[H]
\centering
\includegraphics[width=0.96\textwidth]{robot-brain-dual-system.png}
\caption{Estructura de sistema dual del cerebro del robot: el VLM descompone las instrucciones y el VLA produce las acciones.}
\end{figure}

\begin{knowledgebox}{Lectura de la figura: por qué separar el VLM y el VLA}
El VLM es adecuado para interpretar instrucciones ambiguas, descomponerlas lógicamente y planificar tareas de carácter más general; el VLA debe producir acciones con baja latencia y, en muchos casos, tiene que ejecutarse en el dispositivo. En los escenarios industriales, el conjunto de tareas es limitado y puede invocarse directamente la interfaz de lenguaje del VLA; en escenarios generales, como el hogar, se necesita más que el VLM participe en la descomposición de tareas complejas.
\end{knowledgebox}

\subsection{Latencia en el dispositivo e inferencia en la nube}

Gao Jiyang señala en particular que, si el modelo que ejecuta las acciones se coloca en la nube, la latencia se vuelve un problema muy difícil de resolver. Por eso el VLA, como modelo de acción, suele tener que desplegarse en el dispositivo. Que el VLM deba participar o no en tiempo real depende de la complejidad del escenario. En escenarios industriales y comerciales con veinte o treinta acciones fijas, quizá no sea necesario invocar el VLM en cada paso; en escenarios domésticos más generales, el VLM sí es un componente indispensable.

\begin{warningbox}{No hay que imaginar el «cerebro del robot» como un único modelo grande}
El cerebro del robot no es un LLM monolítico. Incluye, como mínimo, descomposición de tareas, generación de acciones, control en el dispositivo, entrada de percepción, retroalimentación de vuelta al sistema y restricciones de seguridad. Reducirlo a «conectar un modelo grande de lenguaje» subestima los problemas de despliegue y de latencia.
\end{warningbox}

\subsection{Ventajas de las empresas emergentes frente a las grandes empresas}

La sección anterior separó el VLM y el VLA; esta responde a la pregunta de la competencia: si las grandes empresas tienen modelos fundacionales, capacidad de cómputo y talento, ¿con qué puede una empresa emergente construir el cerebro del robot? La respuesta de Gao Jiyang no es afirmar de forma general que son «más rápidas», sino volver al punto de entrada de los datos.

Gao Jiyang divide los factores de éxito para construir un buen VLA en datos, algoritmos, capacidad de cómputo/infraestructura y talento. Las grandes empresas son fuertes en capacidad de cómputo, infraestructura y talento, pero con frecuencia carecen de datos reales de manipulación; las empresas emergentes que tienen capacidad de fabricar el robot completo pueden ser más fuertes en el punto de entrada de los datos. Sobre todo en China, la cadena de suministro y la capacidad de iterar el hardware facilitan que las empresas emergentes integren el robot completo, los escenarios y la recolección de datos.

\subsection{Resumen de la sección}

Lo decisivo en el cerebro del robot no es solo «quién tiene el modelo grande más potente», sino si se logra contar con un VLA ejecutable en el dispositivo, una descomposición adecuada a cargo del VLM en el nivel superior, un ciclo cerrado de datos reales y la capacidad de desplegar el robot completo.
\section{La salida de Xu Huazhe (许华哲): innovación algorítmica y barreras en la cadena de valor}

Después del cerebro del robot, la entrevista llega a una pregunta delicada pero de gran valor: ¿qué revela la salida de un cofundador? Esta sección no emite juicios sobre rumores; la trata como una ventana para entender cómo se jerarquizan las barreras de entrada de una empresa de robótica.

Cuando se grabó la entrevista, Xu Huazhe, cofundador de Xinghaitu (星海图), estaba por dejar la empresa. El entrevistador insistió: ¿significa esto que, en la etapa actual, la innovación algorítmica no es importante para las empresas de robótica? La respuesta de Gao Jiyang (高继扬) fue clara: la innovación algorítmica es importante, pero no puede existir de forma independiente, separada de toda la cadena de valor de la inteligencia corporizada.

\begin{figure}[H]
\centering
\includegraphics[width=0.96\textwidth]{embodied-ai-value-chain.png}
\caption{Cadena de valor de la inteligencia corporizada y ciclos de difusión: robot completo, clientes, datos, AI infra, algoritmos/modelos.}
\end{figure}

\begin{knowledgebox}{Lectura de la figura: el ciclo de difusión determina las barreras de corto plazo}
Los ciclos de difusión de la figura provienen de las estimaciones hechas en la entrevista: el robot completo y la cadena de suministro requieren entre 12--18 meses; los canales de clientes, al menos 6 meses; el sistema de datos, de 6--12 meses adicionales a partir del robot completo; en cambio, un equipo de primer nivel puede alcanzar un artículo de algoritmos o un método de código abierto en apenas 2--3 meses. Por ello, las barreras de corto plazo no dependen solo del algoritmo, sino de toda la cadena.
\end{knowledgebox}

\subsection{Por qué el algoritmo sigue siendo importante}

Antes se usaron los ciclos de difusión para mostrar que el algoritmo no es necesariamente la barrera más sólida, pero eso no significa que pueda ignorarse. Una lectura más precisa es esta: el algoritmo debe integrarse en activos de ciclo más largo para que la innovación se convierta en barrera.

Gao Jiyang no menospreció el algoritmo. Subrayó que el equipo de algoritmos de Xinghaitu es muy capaz y que seguirá innovando. Pero en el entorno actual, en el que el código abierto y los artículos se difunden muy rápido, el algoritmo por sí mismo es más fácil de replicar; sin el respaldo del robot completo, los datos, la infra, los canales y el valor para el cliente, es difícil que la innovación algorítmica se convierta en una defensa de largo plazo.

\begin{importantbox}{Innovación pragmática}
«Innovación pragmática» no significa dejar de innovar, sino evaluar primero el ROI de la innovación dentro de toda la cadena de valor: cuánto aporta a la estrategia de largo plazo, cuánto aporta a los resultados de corto plazo y si el robot completo, los datos y el ciclo cerrado con los clientes pueden amplificarla.
\end{importantbox}

\subsection{Valores y decisiones organizacionales}

Una vez tratados el algoritmo y las barreras, la pregunta pasa de manera natural a la organización: cuando las personas, la dirección y los recursos ya no coinciden por completo, ¿cómo decide el equipo fundador? Ahí es donde la «innovación pragmática» se pone realmente a prueba.

Gao Jiyang también explicó la salida de Xu Huazhe como un cambio de etapa de la organización: en cada etapa, la empresa necesita distintos tipos de personas para crear valor, y también necesita hacer ajustes con apego a los hechos. Definió los valores como dos tipos de decisiones: ante una elección de dirección, qué se elige y qué no; ante el reparto de beneficios, a quién se le asigna y a quién no. Esto se acerca más a la lógica real de una organización emprendedora que preguntarse «quién tiene la razón y quién no».

\begin{warningbox}{No atribuir la salida de un cofundador a una sola causa}
Los cambios en el equipo fundador pueden deberse a una combinación de dirección, etapa, responsabilidades, valores y la voluntad personal de emprender. Para un observador externo, lo más importante es ver si la organización, tras el ajuste, sigue produciendo resultados, y no fijarse solo en el cambio de personal en sí.
\end{warningbox}

\subsection{Resumen de la sección}

El pasaje sobre la salida de Xu Huazhe, en realidad, deja claras las barreras de una empresa de inteligencia corporizada: el algoritmo es importante, pero debe integrarse en la cadena del robot completo, la cadena de suministro, los datos, la infra, los modelos, la distribución y el valor para el cliente.
\section{Pares, financiamiento y complejidad organizacional}

Los capítulos anteriores trataron la cadena técnica y de producto; este capítulo añade las coordenadas externas de la empresa emergente. Gao Jiyang (高继扬) mencionó varias veces a sus pares durante la entrevista: aprender de Yushu (宇树) en el robot completo y la cadena de suministro, aprender de Physical Intelligence en modelos fundacionales y densidad de talento, y observar también lo que hacen empresas como Zhiyuan (智元) en gestión, propiedad intelectual y organización. Esta forma de expresarlo resulta interesante: no ve a sus pares solo como competidores, sino que trata a la industria como un sistema del que se puede aprender.

\subsection{Qué aprender de los pares}

Esta subsección desglosa «aprender de los pares» en fuentes de capacidad más concretas: las empresas de hardware ofrecen un modelo de cadena de suministro, las empresas de modelos ofrecen un modelo de inteligencia y los equipos directivos maduros ofrecen un modelo de organización. Vistos así, los competidores también son una fuente de conocimiento de la industria.

En el caso de Yushu, le interesa su profundidad en la cadena de suministro: capacidades de base como engranajes, motores, carcasas y simulación electromagnética. En el caso de Physical Intelligence, le interesan su posición de liderazgo en la dirección de la inteligencia, su densidad de talento y su densidad de capital. En el caso de Zhiyuan, le interesa cómo un equipo directivo maduro administra una empresa de inteligencia corporizada, incluidas la propiedad intelectual, los ajustes organizacionales y las decisiones de gestión pragmáticas.

\begin{knowledgebox}{Las tres dimensiones del aprendizaje de los pares}
\begin{tabularx}{\textwidth}{p{0.22\textwidth}p{0.32\textwidth}X}
\toprule
Referente & Qué se puede aprender & Relevancia para Xinghaitu (星海图) \\
\midrule
Yushu & Robot completo, cadena de suministro, desarrollo propio de componentes y profundidad de manufactura & Completa la base física de una empresa de robótica. \\
Physical Intelligence & Modelos fundacionales, densidad de talento e inversión en algoritmos de frontera & Referente para la capacidad aguas arriba del cerebro del robot. \\
Zhiyuan & Equipo directivo, propiedad intelectual, ajustes organizacionales y decisiones de operación & Aprender cómo una organización compleja entrega resultados de forma sostenida. \\
\bottomrule
\end{tabularx}
\end{knowledgebox}

\subsection{El aumento de la valuación no es la raíz de los problemas organizacionales}

En la entrevista se menciona que la valuación de Xinghaitu aumentó considerablemente en dos años y que el equipo pasó de poco más de diez personas a más de doscientas. Gao Jiyang sostiene que lo que realmente provoca los problemas organizacionales no es el aumento de la valuación en sí, sino que la organización se vuelve más compleja, que el scope de las tareas se amplía con rapidez, si los miembros originales y el equipo fundador pueden seguir el ritmo de crecimiento, y si es posible incorporar a tiempo a personas con más experiencia.

Las empresas de inteligencia corporizada enfrentan además un problema organizacional particular: la cadena de suministro del robot completo exige procesos, disciplina y calidad; el equipo de inteligencia exige densidad de talento, innovación y ensayo y error rápido. Estos dos domain tienen lenguajes de gestión distintos, ritmos distintos y formas distintas de tolerar errores. Reunirlos en una misma empresa es un desafío organizacional inherente al emprendimiento en inteligencia corporizada.

\begin{importantbox}{La doble naturaleza organizacional de las empresas de inteligencia corporizada}
De un lado está el hardware y la cadena de suministro, que requieren disciplina, procesos, calidad y control de costos; del otro, la inteligencia y los modelos, que requieren talento de alta densidad, exploración e iteración rápida. Poder gestionar ambas culturas al mismo tiempo es la cuestión central para que una empresa de robótica pueda crecer.
\end{importantbox}

\subsection{Resumen de la sección}

Xinghaitu no compite en el vacío. Por un lado aprende de sus pares; por el otro, atiende el financiamiento, la expansión organizacional y la gestión de dos domain. Para una empresa de inteligencia corporizada, que una ruta tecnológica llegue a cumplirse depende, en última instancia, de que la organización pueda soportar la complejidad.
\section{De la demo a la productividad: un marco de aceptación}

En este episodio aparece una y otra vez una pregunta de evaluación implícita: ¿en qué momento un robot puede considerarse utilizable? Gao Jiyang (高继扬) describe las etapas de capacidad desde la demo en video, la demo en oficina y la demo en entornos reales (\emph{wild demo}) hasta los escenarios de productividad. Esta división es más útil que fijarse solo en las métricas del modelo, porque lo que el robot termina por enfrentar no es una tabla de clasificación, sino clientes reales y tareas reales.

\subsection{Cuatro niveles de aceptación}

Aquí se organizan las formas de demo mencionadas en la entrevista en una tabla de aceptación de cuatro niveles. El propósito de la tabla no es inventar términos nuevos, sino recordar al lector que la capacidad de un robot debe pasar de «verlo tener éxito una vez» a «generar valor de forma continua para el cliente».

\begin{longtable}{p{0.22\textwidth}p{0.34\textwidth}p{0.35\textwidth}}
\toprule
Etapa & Qué demuestra & Qué no demuestra \\
\midrule
Demo in Video & Es visualmente atractiva; demuestra que una acción puede captarse en video & No demuestra estabilidad, reproducibilidad ni adaptación al sitio. \\
Demo in Office & Puede presentarse en vivo en el entorno controlado de la empresa & No demuestra que funcione en distintos lugares, con distintos objetos ni en los procesos de distintos clientes. \\
Demo in the Wild & Puede desplegarse y presentarse con clientes, en ferias o en escenarios en otros países & No equivale a operación prolongada sin supervisión ni a un bajo costo de mantenimiento. \\
Escenario de productividad & El cliente está dispuesto a usarlo de forma continua y a pagar por el valor que recibe & Aún requiere validar la entrega a escala, el mantenimiento y el ciclo cerrado de datos. \\
\bottomrule
\end{longtable}

\begin{warningbox}{El sesgo en la percepción de los robots que generan los videos}
La idea que el público general tiene de los robots suele provenir de videos cortos. Los videos amplifican los fragmentos exitosos y minimizan el número de fallas, el tiempo de preparación, las limitaciones del escenario y el costo de mantenimiento. Por ello, para evaluar a una empresa de robótica no basta con ver las demos: hay que observar los ciclos, los envíos, la recompra de los clientes, el mantenimiento y el retorno de datos.
\end{warningbox}

\subsection{El sitio del cliente es la evaluación definitiva}

Gao Jiyang dice que, si un cliente opina que el producto no es bueno, él acude de inmediato al sitio del cliente para resolverlo y convierte cada problema individual en la solución sistemática de toda una clase de problemas. Esta práctica corresponde a la forma real de evaluar la inteligencia corporizada: no es una calificación única en una tabla de clasificación fuera de línea, sino exponer problemas una y otra vez en el sitio del cliente, resolverlos, consolidar los procesos y mejorar el producto.

\begin{importantbox}{La triple validación de los escenarios de productividad}
Para que una capacidad robótica entre en un escenario de productividad debe superar, como mínimo, una triple validación: que el cliente esté dispuesto a pagar, que el sistema complete las tareas de forma estable y que el sistema de mantenimiento y el ciclo cerrado de datos absorban los problemas de manera continua.
\end{importantbox}

\subsection{Resumen de la sección}

Para pasar de la demo a la productividad, una empresa de robótica debe superar cuatro umbrales: que el sistema sea reproducible, desplegable, mantenible y que el cliente pague por él. La verdadera capacidad del modelo tiene que redefinirse a la luz de estos umbrales.
\section{«Bajar a la tierra»: la forma real del idealismo}

En la segunda mitad de la entrevista, el conductor insiste en la vision técnica. La visión que plantea Gao Jiyang (高继扬) es entrenar a los robots igual que se capacita a un empleado: con unas cuantas demostraciones y algunas prácticas por cuenta propia, que el robot complete tareas de forma estable y autónoma en su entorno. Para lograr esa experiencia, Xinghaitu (星海图) necesita tres componentes: el modelo base, las herramientas de postentrenamiento y el robot completo.

\begin{figure}[H]
\centering
\includegraphics[width=0.94\textwidth]{go-to-soil-robotics.png}
\caption{«Bajar a la tierra»: por qué emprender en robótica no es romántico.}
\end{figure}

\begin{knowledgebox}{Lectura de la figura: no ser romántico no significa carecer de idealismo}
A la izquierda, las aplicaciones de software y de modelos grandes tienden a concentrarse en el modelo, el crecimiento y la difusión; a la derecha, una empresa de robótica tiene que hacerse cargo del robot completo, la cadena de suministro, las instalaciones del cliente, el mantenimiento, la seguridad y el retorno de datos. La «tierra» del centro representa la fricción del mundo real. El idealismo de emprender en robótica no es un eslogan, sino mantener el objetivo después de enfrentar esas fricciones todos los días.
\end{knowledgebox}

\subsection{Cómo reconocer las promesas vacías}

Gao Jiyang considera que no es posible reconocer en un instante si una empresa de robótica hace promesas vacías; hay que observar un periodo: qué dijo el equipo hace uno o dos años, qué hizo durante el último año y si cumple dentro de un año. En la industria robótica, describir el futuro es necesario, porque los empleados, los inversionistas, los proveedores y los clientes necesitan creer en él; lo decisivo es si el equipo convierte de manera constante esa descripción del futuro en realidad.

\begin{importantbox}{La diferencia entre las promesas vacías y la narrativa estratégica}
Describir el futuro no es el problema; el problema es si ese futuro se desglosa en resultados por etapas, asignación de recursos, valor para el cliente y entregas verificables. Una narrativa que se cumple es estrategia; solo la que no se cumple es una promesa vacía.
\end{importantbox}

\subsection{El lobo: empuje y tenacidad}

Cuando le preguntan a qué animal se parece Xinghaitu, Gao Jiyang responde que lo más cercano es el lobo, aunque subraya que todas las empresas tienen algo de lobo. Aquí el «lobo» no alude a una cultura brusca, sino al empuje, la tenacidad y la determinación para alcanzar metas frente a las entregas a clientes y las fechas de lanzamiento. La intensidad de la competencia en la industria de la inteligencia corporizada exige que el equipo tenga idealismo y, al mismo tiempo, calcule con pragmatismo, todos los días, la relación entre inversión y resultados, y el valor a largo plazo.

\subsection{Resumen de la sección}

«Bajar a la tierra» es la metáfora industrial más importante de este episodio. Una empresa de robótica debe encontrar el equilibrio entre idealismo y pragmatismo: la visión tiene que llegar lejos, pero cada día hay que resolver problemas de hardware, cadena de suministro, clientes, datos y organización.
\section{Términos clave: índice de palabras clave de este episodio}

Las secciones anteriores ya cubrieron lo personal, la organización, la tecnología y el modelo de negocio. Aquí se reúnen los términos en un solo lugar, para facilitar las referencias cruzadas con EP121, EP109, EP098 y otras notas relacionadas con robots.

\begin{longtable}{p{0.24\textwidth}p{0.32\textwidth}p{0.35\textwidth}}
\toprule
Término & Explicación en una frase & Papel en este episodio \\
\midrule
Physical World AI & IA que actúa directamente sobre tareas físicas reales & Base común de la conducción autónoma y la inteligencia corporizada. \\
Trayectoria de Waymo & robotaxi, ingeniería de sistemas, ejecución a largo plazo & Aporta formación en descomposición de sistemas y en ingeniería. \\
Trayectoria de Momenta & entrega en producción en serie, valor para el cliente, ciclo de datos & Aporta formación en conversión en producto y en organización. \\
Equipo completo & cuerpo del robot, hardware y punto de entrada del cómputo en el dispositivo & Base de la recolección de datos y de la entrega al cliente. \\
Cadena de suministro & sistema de componentes, manufactura, calidad y entrega & Su ciclo de difusión es largo; es una fuente de ventaja defensiva. \\
Data Recipe & receta de recolección, limpieza, entrenamiento, evaluación y retorno de datos & Determina cómo los datos reales se convierten en capacidades. \\
VLM & Vision-Language Model, para descomponer instrucciones y razonar en el nivel superior & Sistema de nivel superior del cerebro del robot. \\
VLA & Vision-Language-Action, genera acciones a partir de visión y lenguaje & Modelo fundacional de acción del cerebro del robot. \\
Demo in the Wild & mostrar capacidades en escenarios reales y abiertos & Más cercana a la generalización real que una demo en video. \\
Ciclo de difusión & tiempo que necesita un competidor para copiar una capacidad & Indicador para juzgar el grosor de la ventaja defensiva. \\
Innovación pragmática & innovar con el valor para el cliente y el ROI como restricción & Uno de los valores organizacionales de Xinghaitu (星海图). \\
\bottomrule
\end{longtable}

\subsection{Resumen de la sección}

Todos estos términos apuntan a una misma conclusión: la competencia en inteligencia corporizada no es una competencia de algoritmos aislados, sino una competencia de sistemas a lo largo de una cadena extensa. El modelo, el equipo completo, la cadena de suministro, los datos, los clientes y la organización deben considerarse en conjunto.
\section{Síntesis y ampliación}

\subsection{Conclusiones principales}

\begin{enumerate}
\item La metodología de Gao Jiyang (高继扬) consiste en inducir y generalizar, razonar hacia atrás a partir del objetivo y gestionar probabilidades; aplica este método de manera constante en los exámenes, los artículos, la ingeniería y el emprendimiento.
\item Waymo forma en ingeniería de sistemas: descomponer, medir, infra y ejecución a largo plazo; pero queda lejos de los clientes y de la operación del negocio.
\item Momenta forma en la entrega para producción en serie: del demo al producto, del research lab al valor para el cliente; la organización tiene que curtirse en proyectos reales.
\item En el modelo de negocio de la conducción autónoma, lo decisivo son los datos y los clientes: quien posee los vehículos, los usuarios y el fondo de utilidades forma con mayor facilidad un ciclo cerrado de datos.
\item Xinghaitu (星海图) eligió construir el robot completo porque el robot completo es la puerta de entrada a los datos y al valor para el cliente, no porque el hardware en sí sea más romántico.
\item La Data Recipe es el mecanismo central con el que una empresa de robótica convierte datos del mundo real en capacidades del modelo.
\item El cerebro del robot no es un único LLM, sino un sistema doble formado por la descomposición de tareas en la capa superior con un VLM y un modelo de acciones VLA, sujeto a las restricciones de latencia en el dispositivo.
\item La innovación algorítmica es importante, pero en la era del código abierto su ciclo de difusión es corto; es más probable que las barreras de largo plazo provengan del robot completo, la cadena de suministro, los datos, los clientes y la organización.
\item «Bajar a la tierra» indica que el idealismo del emprendimiento en robótica debe tomar la forma de pragmatismo, ROI, entregas y presencia en las instalaciones del cliente.
\end{enumerate}

\subsection{Preguntas abiertas}

Al final se dejan preguntas abiertas porque la inteligencia corporizada todavía está en una etapa en la que las rutas no han convergido. Este episodio ofrece un marco realista contundente, pero aún hay que seguir observando cómo se materializará ese marco en los productos y modelos de los próximos años.

\begin{itemize}
\item ¿Una empresa de inteligencia corporizada debe atender primero el mercado de desarrolladores o entrar directamente en escenarios productivos?
\item ¿La generalidad de los VLA provendrá de la escala del modelo, de la escala de los datos o del ciclo cerrado entre el robot completo y el escenario?
\item ¿Puede una empresa emergente con capacidad de construir el robot completo formar una ventaja de largo plazo en datos frente a las grandes empresas?
\item A medida que el ciclo de difusión de los algoritmos se acorta, ¿el foso defensivo de las empresas de robótica se desplazará cada vez más hacia la cadena de suministro, los clientes y la recipe de datos?
\item ¿Puede la ventaja de la cadena de suministro de China traducirse de forma sistemática en una ventaja en modelos base de inteligencia corporizada?
\end{itemize}

\subsection{Lecturas adicionales}

\begin{itemize}
\item EP121, entrevista con Tan Jie (谭捷) de DeepMind: robots, transferencia entre encarnaciones, modelos del mundo y Gemini Robotics.
\item EP109, conversación con Xie Chen (谢晨) sobre la escasez de datos en robótica: simulación, datos sintéticos y la industria de los datos.
\item EP098, explicación de artículos clásicos sobre modelos base para robots y VLA: para entender la evolución técnica de los VLA.
\item Materiales sobre arquitectura de sistemas de conducción autónoma, conducción autónoma de extremo a extremo y el modelo de negocio del robotaxi.
\item Materiales sobre la cadena de suministro de robots, el despliegue en el dispositivo, los datos de teleoperación y las herramientas de postentrenamiento.
\end{itemize}

\begin{importantbox}{El juicio final}
El juicio que más vale la pena llevarse de este episodio es este: la inteligencia corporizada no pasa de manera natural de «el modelo es muy potente» a «el robot es utilizable», sino que forma un ciclo cerrado en el que intervienen conjuntamente el robot completo, los clientes, los datos, el modelo y la organización. El verdadero romanticismo técnico quizá no consista en hablar del futuro desde las nubes, sino en estar dispuesto a hundir la cabeza en la tierra y construir el futuro poco a poco.
\end{importantbox}

\end{document}
